\section{The phenomenon}
% Frames P1--P4.

\subsection{What holds together over a long exchange}
% P1

A language model asked to reason sustains something that reads as a mind. Over an exchange of several
hours it keeps the commitments it made early on, returns to a thought it left unfinished, and answers
the interlocutor's mood in the way the interlocutor would expect of a person. Over a chain of hundreds
of reasoning steps it keeps the goal of the chain available while working on a subgoal, and it notices
when a subgoal has made the goal unreachable. Calling this imitation is easy. Saying how an imitation
could be that stable is hard, and that is the question this paper starts from.

\subsection{What the deflationary answer predicts}
% P2

Consider what follows if a language model reproduces the surface statistics of text with nothing behind
them. Then it should fail the way low-order Markov processes fail. A thread introduced in the third
paragraph should be dropped by the sixth. An attitude toward the interlocutor should flip when the
topic changes, because nothing carries it. A motive attributed to a character should not survive the
page on which it was attributed. These failures are cheap: the space of continuations that produce them
is vastly larger than the space that avoids them, and a long horizon offers many occasions to fall into
it. In current models such failures are rare, and their rate falls with scale and training rather than
staying fixed. Something keeps the process out of a large space of cheap failures for a long time. That
something is the object of this paper.

The observation is weak on its own. It does not say that the thing keeping the process on course
resembles anything in us, and the history of the subject is full of inferences from surface competence
to inner life that did not survive. What the observation does is set a requirement on any account: the
account has to name a structure that could hold a long exchange together, and it has to say how much
of that structure a given system has.

\subsection{Nagel's question, put to a model}
% P3

\citet{nagel1974} asked what it is like to be a bat in order to mark a limit. The bat's experience, if
it has any, is organized around echolocation, and no description from outside crosses over into it. Put
the same question to a language model and it loses part of its force and gains another part. It loses
force because the model shares our language, and the bat does not: whatever the model has, it has in a
medium we can read, and it will answer questions about it in our words. It gains force because the
model's substrate is further from ours than any animal's. A bat has neurons, a body, metabolic needs, a
developmental history and a lifetime of continuous sensory contact with a world. A model has a sequence
of tensor operations over a context window.

So we do not ask whether there is something it is like to be a model. We ask what a model has of the
functions that in us go under the name of consciousness, and how much of each it has. The change of
question is not a retreat from the hard case. It is the only form of the question we know how to
answer, and Section~\ref{sec:measurement} shows that in this form it has answers that can come out
either way.

\subsection{The two shortcuts}
% P4

Two answers arrive before any argument is made. The first: there is nothing there, it is a text
predictor, and the impression of a mind is the reader's projection. The second: there is someone there,
it says so, treat it as a person. Neither is a position in the sense that would let one derive a
prediction from it and be wrong. The first cannot say what it would take to change its mind; the second
cannot say what it has attributed. Section~\ref{sec:reader} argues that both are what a mind reaches
for when it lacks the concepts to reach for anything else, and that their popularity is evidence about
the cost of concepts rather than about machines. The paper is about what lies past them.
